%==============================================================================
% System-Heterogeneity section — drop-in LaTeX for Overleaf
% Sibling to thesis_ready/writing/09_overleaf_ready.tex (statistical het).
% Copy this entire file as a separate .tex document in the same Overleaf
% project, OR \input it into your main thesis document after the statistical
% heterogeneity section.
%
% All HPC-pending values use \TODOhpcSH{} markers. Toggle \showtodoshtrue/false
% near the top to render markers visibly (default) or hide them.
%
% Compiles standalone (with \documentclass etc. at top) or as an \input{...}
% chunk (delete the document scaffolding).
%==============================================================================
\documentclass[11pt,a4paper]{article}

\usepackage[utf8]{inputenc}
\usepackage[T1]{fontenc}
\usepackage{graphicx}
\usepackage{booktabs}
\usepackage{amsmath, amssymb}
\usepackage{microtype}
\usepackage{hyperref}
\usepackage{xcolor}
\usepackage[margin=2.5cm]{geometry}
\usepackage{setspace}
\onehalfspacing

\providecommand{\citep}[1]{(#1)}
\providecommand{\citet}[1]{#1}

% --- TODO markers for HPC-pending system-het values -------------------------
% These are deliberately distinct from \TODOhpc in the statistical-het file,
% so the two sections can co-exist in the same Overleaf project without
% macro collisions if you choose to merge them.
\newif\ifshowtodosh
\showtodoshtrue
\newcommand{\TODOhpcSH}[1]{%
  \ifshowtodosh\textcolor{red}{\textbf{[SH-HPC: #1]}}\else#1\fi}

\title{System Heterogeneity: FedProx under Variable Client Compute Budgets}
\author{}
\date{}

\begin{document}
\maketitle

%==============================================================================
\section{Motivation and Background}
%==============================================================================

The statistical-heterogeneity experiment of the preceding section
demonstrated that FedProx \citep{li2020fedprox} outperforms FedAvg
\citep{mcmahan2017fedavg} on a deliberately non-IID partition of
DermaMNIST. That experiment isolated the data-distribution axis of
federated heterogeneity. The present section isolates the second canonical
axis identified by the original FedProx paper: \textbf{system
heterogeneity} -- variability in the local computational budget each
client can contribute per communication round.

\paragraph{Terminology.} Throughout this section we use ``system
heterogeneity'' as shorthand for what is strictly
\emph{computational-work heterogeneity}: variability in the number of
local SGD steps each client performs per round. We do not simulate
hardware speed, network latency, or device availability; the
manipulation acts on the algorithm via heterogeneous $\tau_i$ only.
Following \citet{li2020fedprox}'s terminology is convenient but
should not be read as a hardware claim.

\paragraph{What FedProx's proximal regulariser was designed for.}
System heterogeneity is, in fact, the \emph{primary} motivation for
FedProx's proximal regulariser. \citet{li2020fedprox} proved (their
Theorem~4) a convergence guarantee for FedProx under $\gamma$-inexact
local updates: each client may submit a parameter vector $w_i^{t+1}$ that
does not fully minimise its local objective, provided the residual
inexactness is bounded. The proximal term $\frac{\mu}{2}\lVert w - w^t \rVert^2$
constrains how far an inexact local update can drift from the global
anchor $w^t$, making partial work from stragglers usable for aggregation
without destabilising the global model. FedAvg has no such guarantee, and
its size-weighted aggregation rule treats partial-work and full-work
updates symmetrically --- which biases the global model toward whichever
clients drifted furthest, regardless of how much work they actually
completed.

\paragraph{One key implementation deviation from Li et al.\ (2020).}
In our implementation, both FedAvg and FedProx aggregate over
\emph{the same} client set every round, including stragglers, weighted
by full dataset size $n_i$. This differs from \citet{li2020fedprox}
§5.2, where the canonical FedAvg straggler protocol \emph{drops}
straggler updates entirely while FedProx \emph{includes} them. Our
design therefore isolates the effect of the proximal term under
heterogeneous local computation without confounding it with
differential client participation. A consequence is that the original
``FedProx tolerates dropping where FedAvg does not'' result of Li
et al.\ (2020) is not directly tested here; what is tested is whether
FedProx's headline statistical-heterogeneity advantage persists or
widens when both algorithms see partial-work updates from a random
subset of clients each round.

\subsection*{Why this complements the statistical-heterogeneity result}

The preceding section answered: \emph{does FedProx help when clients have
heterogeneous data distributions?} This section answers a logically
distinct question: \emph{does FedProx help when clients have heterogeneous
computational budgets?} Both are positive answers in the original Li et
al.\ paper, but on different benchmarks (Synthetic, FEMNIST, Shakespeare,
Sentiment140) -- not on medical imaging. Demonstrating both effects on
DermaMNIST under the same partition, same paired-seed protocol, and same
statistical machinery gives a coherent two-axis defence of FedProx's
theoretical motivation.

\subsection*{Scope: FedProx vs.\ FedNova for heterogeneous local steps}

This section restricts itself to FedAvg vs.\ FedProx, matching the
project brief and the protocol of the preceding statistical-heterogeneity
section. We acknowledge that for the specific case of \emph{heterogeneous
numbers of local update steps}, the more specialised competitor algorithm
is \textbf{FedNova} \citep{wang2020fednova}, which argues that
heterogeneous local steps create objective inconsistency under naïve
weighted aggregation and proposes a normalised-averaging rule to address
it. FedNova's mechanism is therefore directly engineered for the regime
we test in this section, while FedProx's proximal-anchor mechanism is
more general (it addresses both client-distribution heterogeneity and
straggler inexactness). A direct FedNova comparison is identified as
the most natural extension of this section and is recommended for
follow-up work; we do not implement it here.

%==============================================================================
\section{Experimental Setting}
%==============================================================================

\subsection{Design choices held identical to the statistical-heterogeneity sweep}

To isolate the system-heterogeneity contribution, every other experimental
parameter is held identical to the headline statistical-heterogeneity
sweep of the preceding section: same \texttt{balanced\_paired\_7\_clients}
partition, same 10 paired seeds $\{42, 123, 456, 789, 999, 2024,$
$31337, 8675309, 161803, 271828\}$, same \texttt{DermMNISTCNN}
architecture, same hyperparameters ($R = 150$, $\mu = 0.01$, batch size
32, SGD with $\eta = 0.01$, momentum $0.9$), and same evaluation
protocol (test set evaluated once per run at the best-val checkpoint).
The only variable is the per-client per-round local-epoch budget.

\subsection{Three conditions}

\paragraph{C0 (baseline, no system heterogeneity).} Every client performs
$E = 20$ local epochs every round. The \emph{design} is identical to
the headline statistical-heterogeneity sweep (same partition, same 10
paired seeds, same hyperparameters), but C0 is re-executed in the
Flower runtime via \texttt{submit\_flower\_C0\_baseline.sh} so that
the H2 contrast (Δ$_c -$ Δ$_{C0}$) consumes runtime-matched data
rather than mixing the pure-PyTorch headline against Flower stragglers.
As a cross-runtime reference, the pure-PyTorch headline produced
$M_\text{FedAvg} = 0.481 \pm 0.025$, $M_\text{FedProx} = 0.508 \pm 0.014$,
$\Delta = +0.027$, paired Wilcoxon $p = 0.020$, rank-biserial
$r_{rb} = +0.818$; the Flower C0 numbers (queued, \TODOhpcSH{value once
sweep lands}) are expected to fall within the documented $\sim 0.005$
runtime noise floor of these values.

\paragraph{C1 (fixed stragglers --- a specialist + a generalist).} Two
designated clients (C5 and C6 in the partition) always perform
$E_\text{straggler} = 5$ local epochs; the remaining five clients perform
$E = 20$. This models a deployment where specific institutions
consistently have weaker compute infrastructure.

\emph{Important caveat on the C1 design.} C5 belongs to the
melanoma/vascular specialty pair (the per-specialty pair that drives the
headline FedProx advantage; see preceding section), and C6 is the
nevi-only generalist. C1 therefore conflates ``system heterogeneity'' with
``a particular clinical role configuration'' --- specifically, the
combination of a high-Δ specialist client and the generalist client both
running slow. Any C1 result must be interpreted as \emph{this particular
configuration} rather than a clean isolation of the system-heterogeneity
effect. The cleanly-isolating test is C2 below, which approximates
the \citet{li2020fedprox} §5.2 random-straggler protocol up to the
aggregation deviation noted above. C1 is retained as an
\textbf{exploratory descriptive cross-check} because of its
operational realism (in real deployments, certain client identities
have predictable hardware constraints), but the H2 (between-condition)
inference in this section relies on C2 as the \textbf{primary
condition}; C1 results are presented only descriptively.

\paragraph{C2 (random stragglers, \citet{li2020fedprox} §5.2-inspired
protocol). PRIMARY condition for inference.} Each communication round,
\textbf{4 of 7 clients per round (approximately 57\%)} are randomly
designated stragglers (sampled without replacement; the fraction is
$\mathrm{round}(0.5 \times 7) = 4$, slightly above 50\% due to integer
rounding). Each straggler performs $E_i \sim
\text{Uniform}\{1, 2, \ldots, 19\}$ local epochs; non-stragglers perform
$E = 20$. The straggler schedule is seeded with the same seed as the
rest of the run, so paired FedAvg/FedProx runs share an identical
schedule and the within-pair $\Delta$ reflects only the algorithmic
difference. This is inspired by the \citet{li2020fedprox} §5.2
random-straggler protocol with two implementation deviations: (i) the
canonical setup includes $E_i = 0$ in the support; we forbid 0 epochs
so every straggler does at least one optimiser step; (ii) the canonical
FedAvg arm drops straggler updates while ours includes them (see the
implementation-deviation paragraph above). C2 therefore tests the
proximal-term effect under heterogeneous partial-work updates, not the
straggler-dropping mechanism that distinguishes Li et al.'s FedAvg
from their FedProx.

\subsection{Pre-registered hypotheses}

\paragraph{H1 (within-condition).} For each system-heterogeneity
condition $c \in \{C1, C2\}$, FedProx outperforms FedAvg on test macro-F1
under paired Wilcoxon ($\alpha = 0.05$, two-sided, $n = 10$).

\paragraph{H2 (between-condition).} The within-pair $\Delta$ is larger
under system heterogeneity than at baseline:
$\mathbb{E}[\Delta_s^{C1}] > \mathbb{E}[\Delta_s^{C0}]$ and
$\mathbb{E}[\Delta_s^{C2}] > \mathbb{E}[\Delta_s^{C0}]$. This is tested
by paired Wilcoxon on the per-seed difference $\Delta_s^c - \Delta_s^{C0}$
for each $c \in \{C1, C2\}$. A positive significant result supports the
canonical FedProx claim that the proximal regulariser is most beneficial
under inexact local updates. Because H2 comprises two related tests (one
per condition), we report both raw and Bonferroni-corrected $p$-values
($\alpha_\text{family} = 0.05$, corrected per-test threshold
$\alpha / 2 = 0.025$). Because C1 is confounded with class composition,
the \textbf{primary interpretation rests on C2} (the Li et
al.\ 2020-inspired protocol, modulo the implementation deviations noted
above) and is robust to whether C1 also reaches significance; C1's
H2 result is reported descriptively, not as a confirmatory test.

%==============================================================================
\section{Results}
%==============================================================================

\subsection{Headline test performance under system heterogeneity}

\begin{table}[h]
\centering
\caption{Test macro-F1 under three conditions of system heterogeneity,
10 paired seeds each. The C0 row is the Flower-runtime baseline
produced by \texttt{submit\_flower\_C0\_baseline.sh}; numbers are
expected within $\sim 0.005$ of the pure-PyTorch headline reference
($0.481 / 0.508 / \Delta\!=\!+0.027 / p\!=\!0.020$). The within-condition
$\Delta$ tests H1; the between-condition $\Delta - \Delta_{C0}$ tests
H2 with Bonferroni correction across active conditions.}
\label{tab:sh-headline}
\small
\begin{tabular}{lccccc}
\toprule
Condition & FedAvg & FedProx & $\Delta$ & Wilcoxon $p$ (H1) & H2 test $p$ (Bonf.) \\
\midrule
C0 (no system het, baseline, Flower) & \TODOhpcSH{FA mean$\pm$SD} & \TODOhpcSH{FP mean$\pm$SD} & \TODOhpcSH{$\Delta$} & \TODOhpcSH{p} & (reference) \\
C1 (fixed stragglers, descriptive only)  & \TODOhpcSH{FA mean$\pm$SD} & \TODOhpcSH{FP mean$\pm$SD} & \TODOhpcSH{$\Delta$} & \TODOhpcSH{p} & \TODOhpcSH{p} \\
C2 (random stragglers, primary)          & \TODOhpcSH{FA mean$\pm$SD} & \TODOhpcSH{FP mean$\pm$SD} & \TODOhpcSH{$\Delta$} & \TODOhpcSH{p} & \TODOhpcSH{p} \\
\bottomrule
\end{tabular}
\end{table}

\subsection{Straggler-tolerance ratio}

To quantify how much each algorithm's performance degrades under system
heterogeneity relative to the no-system-het baseline, we define the
\emph{straggler-tolerance ratio} for algorithm $a$ in condition $c$:
\[
\rho_a^c = \frac{M_a^c}{M_a^{C0}},
\]
where $M$ is mean test macro-F1. Values close to 1 indicate the algorithm
preserves its no-system-het performance; lower values indicate worse
tolerance to stragglers.

\begin{table}[h]
\centering
\caption{Straggler-tolerance ratios. Higher = more tolerant. The prediction
$\rho_\text{FedProx}^c > \rho_\text{FedAvg}^c$ supports the original
FedProx motivation.}
\label{tab:sh-tolerance}
\small
\begin{tabular}{lcc}
\toprule
Condition & $\rho_\text{FedAvg}$ & $\rho_\text{FedProx}$ \\
\midrule
C0 (baseline) & $1.000$ & $1.000$ \\
C1 (fixed stragglers)  & \TODOhpcSH{$\rho_{FA}^{C1}$} & \TODOhpcSH{$\rho_{FP}^{C1}$} \\
C2 (random stragglers) & \TODOhpcSH{$\rho_{FA}^{C2}$} & \TODOhpcSH{$\rho_{FP}^{C2}$} \\
\bottomrule
\end{tabular}
\end{table}

\subsection{Per-class effects under system heterogeneity}

\TODOhpcSH{This subsection will report per-class test F1 for each condition,
using the same Holm-corrected paired Wilcoxon protocol as the statistical-
heterogeneity section. The expectation is that the minority classes most
sensitive to client drift (melanoma in particular) will be most helped by
FedProx under system heterogeneity, paralleling but exceeding the
statistical-heterogeneity per-class result ($\Delta_\text{melanoma} = +0.114$,
$p_{\text{Holm}} = 0.041$).}

%==============================================================================
\section{Discussion}
%==============================================================================

\TODOhpcSH{This section will discuss whether the H1 and H2 results match
the pre-registered predictions:

\begin{itemize}
\item If H1 confirmed: FedProx's advantage under non-IID data extends to
the system-heterogeneity regime; the proximal regulariser is broadly
useful for federated medical-imaging deployment.
\item If H2 confirmed: the system-heterogeneity advantage is larger than
the statistical-heterogeneity advantage, consistent with FedProx's
original theoretical motivation in \citet{li2020fedprox} that the
proximal term was introduced primarily to handle inexact local updates
from stragglers.
\item If H2 not confirmed but H1 confirmed: FedProx remains useful under
system heterogeneity but the magnitude is similar to the
statistical-heterogeneity regime --- the two effects do not combine
multiplicatively, suggesting the proximal term's primary contribution is
constraining drift regardless of its source.
\item If H1 not confirmed: a substantive negative finding worth
reporting; the statistical-heterogeneity claim is unaffected but the
canonical system-heterogeneity FedProx claim does not transfer to medical
imaging in our setup.
\end{itemize}}

%==============================================================================
\section{Summary of Heterogeneity Experiments}
%==============================================================================

This section closes a two-experiment investigation:

\begin{table}[h]
\centering
\caption{Summary of the two heterogeneity experiments. Each tests a
distinct claim about FedProx and uses the same statistical machinery
(paired Wilcoxon, 10 seeds, balanced\_paired\_7\_clients partition).}
\label{tab:sh-summary}
\small
\begin{tabular}{lll}
\toprule
& Statistical heterogeneity (preceding section) & System heterogeneity (this section) \\
\midrule
What varies & Partition (class allocation) & Per-client local-epoch budget \\
Tests FedProx claim & Drift mitigation under non-IID & $\gamma$-inexact updates from stragglers \\
Headline result & $\Delta = +0.027$, $p = 0.020$ & \TODOhpcSH{pending} \\
Effect size & $r_{rb} = +0.818$ & \TODOhpcSH{pending} \\
\bottomrule
\end{tabular}
\end{table}

%==============================================================================
\section*{References}
%==============================================================================
\begin{description}\setlength{\itemsep}{2pt}
\item Li, T., Sahu, A.~K., Zaheer, M., Sanjabi, M., Talwalkar, A., \&
Smith, V. (2020). Federated optimization in heterogeneous networks. In
\textit{Proceedings of Machine Learning and Systems}, \textit{2},
429--450.
\item McMahan, B., Moore, E., Ramage, D., Hampson, S., \& y~Arcas,
B.~A. (2017). Communication-efficient learning of deep networks from
decentralized data. In \textit{Proceedings of AISTATS} (pp.\ 1273--1282).
\item Wang, J., Liu, Q., Liang, H., Joshi, G., \& Poor, H.~V. (2020).
Tackling the objective inconsistency problem in heterogeneous federated
optimization. \textit{NeurIPS}, 33, 7611--7623.
\item Karimireddy, S.~P., Kale, S., Mohri, M., Reddi, S.~J., Stich, S.~U.,
\& Suresh, A.~T. (2020). SCAFFOLD: Stochastic controlled averaging for
federated learning. In \textit{ICML} (pp.\ 5132--5143).
\end{description}

\end{document}
